\documentclass[10pt]{article}
\usepackage[margin=20mm]{geometry}
\usepackage{amsmath,amssymb,mathtools,fontspec,unicode-math,xurl,hyperref}
\IfFontExistsTF{DejaVu Serif}{\setmainfont{DejaVu Serif}}{\setmainfont{Latin Modern Roman}}
\setmathfont{latinmodern-math.otf}
\setlength{\parindent}{0pt}\setlength{\parskip}{6pt}\sloppy
\title{Technical review of
The total rank conjecture in characteristic 2}
\author{AI-assisted initial scientific review}
\date{}
\begin{document}
\maketitle
Abstract: The characteristic-two length argument is consistent under the explicit finite-length, nonzero-homology hypotheses of Theorem 4.1. The acquired manuscript nevertheless defines the codimension of total homology using a maximum. With that definition its central extension to short complexes is false: a minimal complex over a regular local ring of dimension three has total rank six while satisfying the printed shortness condition. Replacing the maximum by the codimension of the total support excludes this example and explains the intended localization. A separate omission of nonzero homology affects introductory and subsidiary formulations. These findings concern the acquired author manuscript; they do not establish a failure of the finite-length module theorem or of the intended corrected result.\par
\section*{Identity, scope, and notation}
This report concerns Keller VandeBogert and Mark E. Walker, The total rank conjecture in characteristic 2, arXiv:2305.09771v2, 11 April 2024, 19 PDF pages, associated with DOI 10.1215/00127094-2024-0027. Page numbers below are one-based PDF pages and agree with the printed pagination. The publisher version was not available for comparison. The primary PDF SHA256 begins `31af4cce625fcb39`; its complete digest, the acquired TeX digest, and all input checks are recorded in the accompanying provenance.\par
The full manuscript text was read. Detailed checks concentrated on the definition of shortness, the symmetric/exterior-square exact sequences, the equality statement in Theorem 3.2, the asymptotic inequalities and positivity in Theorem 4.1, and the reduction in Theorem 5.1. The background theorems on Dold–Kan, Dutta limits, the New Intersection Theorem, and characteristic-zero transfer are dependencies; this review is not an independent proof of those theorems.\par
Write \(H(P)=\bigoplus_jH_j(P)\). Over a local ring, \(\beta(P)\) denotes the sum of the ranks of a minimal free representative, as in Notation 1.13. A complex with all differential entries in the maximal ideal is already minimal.\par
\section*{A false domain for the short-complex extension}
Problem 1 (Major: maximum codimension in Definition 1.5).  Definition 1.5, p. 2, defines a short complex supported in degrees \(0,\ldots,c\) by \[c=\operatorname{codim}_R(H(P)):=\max_j\{\operatorname{codim}_R H_j(P)\}.\] Conjecture 1.6, p. 3, then asserts the total rank bound \(2^c\); Theorem 5.1, p. 16, asserts this conjecture for all Noetherian algebras over a field. The printed maximum admits the following counterexample.\par
Let \(k=\mathbf F_2\), \(R=k[x,y,z]_{(x,y,z)}\), and take \[P=\bigl(0\longrightarrow R
       \xrightarrow{\,0\,}R
       \xrightarrow{\,0\,}R^3
       \xrightarrow{\,[\,x\ y\ z\,]\,}R
       \longrightarrow0\bigr),\] in degrees \(3,2,1,0\). The ring is a connected, Noetherian, regular local field algebra of dimension three. All four terms are nonzero and free, all compositions vanish, and the complex is minimal. Its homology is \[H_0(P)=k,\qquad
 H_1(P)=\ker[\,x\ y\ z\,],\qquad
 H_2(P)=H_3(P)=R.\] The kernel in degree one is nonzero, since it contains \((y,x,0)\). It is a submodule of a free module over a domain and has zero annihilator. Consequently the four homology codimensions are \(3,0,0,0\), respectively. In particular, \[\max_{0\le j\le3}\operatorname{codim}_RH_j(P)=3,\qquad
 \sum_j\operatorname{rank}_RP_j=\beta(P)=6<8=2^3.\] All homology groups in the displayed interval are nonzero. Here the printed maximum is read over that interval, equivalently over the nonzero homology groups in this example. If all zero groups outside the interval are included with codimension \(+\infty\), Definition 1.5 instead assigns no finite \(c\) to a bounded complex and itself requires clarification. Under the finite interpretation, this nonexact complex meets the printed shortness condition. Neither minimality nor a requirement that the endpoint terms be nonzero removes the counterexample.\par
\subsection*{Dependency and counterevidence}
The first sentence of the proof of Theorem 5.1 localizes at a prime of height \(c\) in the support and claims to obtain a tiny complex. For this example, localization at the maximal ideal leaves the positive-rank homology in degrees \(1,2,3\); these modules do not have finite length. Thus the displayed reduction does not apply to the class defined using the maximum.\par
There is strong evidence for a different intended definition. Page 3 identifies tiny complexes with complexes having finite-length homology, and Section 2.2, p. 10, restates that finite-length, nonzero-homology condition explicitly. The operative Theorem 4.1, p. 13, also requires finite-length, nonzero homology. The counterexample above violates those finite-length hypotheses and therefore does not refute that theorem, nor the original total rank conjecture for finite-length modules.\par
\subsection*{Correction boundary}
For a nonzero total homology module, the codimension of its support is \[\operatorname{codim}_R\!\left(\bigoplus_j H_j(P)\right)
 =\min_{H_j(P)\ne0}\operatorname{codim}_R H_j(P).\] Replacing the maximum by this minimum excludes the example: its total support has codimension zero. Localizing at a minimal prime of that support of height \(c\) then makes every surviving homology module finite length over the local ring of dimension \(c\), which is the reduction used on p. 16. This is a concrete correction of the domain, not a proof that the uncorrected assertion holds. The corrected result still depends on the subsequent tiny-complex and transfer arguments. The Major classification concerns a demonstrated false central statement under its printed definition; it does not assert irreparability.\par
\section*{Nonzero homology in subsidiary formulations}
Problem 2 (Minor: omitted nonzero hypothesis).  Introductory Theorem 1.7, p. 3, and Theorem 4.10, pp. 15–16, omit \(H(P)\ne0\). The latter divides by \(h^\infty(P)\). Notation 1.13, p. 6, permits exact objects in \(\mathrm{Perf}^{\mathrm{fl}}(R)\), so membership in that category does not supply the missing condition.\par
Set \(R=\mathbf F_2[[x,y]]\), \(d=2\), and let \(P\) be the contractible complex \(R\xrightarrow{1}R\) in degrees \(1,0\), extended by zero to the displayed intervals of these theorems. This ring is complete and has perfect residue field. Every Frobenius iterate is again contractible, and \[H(P)=0,\quad \beta(P)=h^\infty(P)=h_1^\infty(P)
       =\chi^\infty(P)=0,\qquad
 \sum_j\operatorname{rank}_RP_j=2.\] The literal inequality in Theorem 1.7 would say \(2\ge4\). The ratio in Theorem 4.10 is \(0/0\), while its subsequent assertion for \(\chi^\infty(P)=0\) would give \(0=\beta(P)\ge2\). The zero complex is an even simpler boundary case.\par
Requiring nonzero homology excludes these examples and supplies the positivity needed in the proof. Theorem 4.1 already states this requirement. The surrounding discussion of total rank conjectures also concerns nontrivial homology, which supports this repair. The issue is classified as Minor because it concerns omitted domain restrictions in these formulations and does not defeat the explicitly restricted central characteristic-two theorem. The analogous nontriviality convention should also be stated when formulating the Section 6 conjecture and applications.\par
\section*{Checks supporting the retained core argument}
\subsection*{Exact sequences and equality}
The two sequences in Lemma 2.16, pp. 9–10, were checked with the characteristic-two Frobenius twist: the map \(x\mapsto x^2\) is linear with that twist. The absolute values in (3.3), p. 12, are present in the page image. They yield the lower bound used in (3.6), p. 13, rather than an unsupported signed equality.\par
For Theorem 3.2 the hypotheses include a nonzero finite-length module of finite projective dimension over a quasi-Roberts ring. The checked inequalities give \[\ell(M)\beta(M)\ge h(P\otimes P)
 \ge 2^d\ell(M)+2\ell(\Lambda^2M).\] If \(\beta(M)=2^d\), then \(\Lambda^2M=0\). Reduction modulo the maximal ideal forces \(M\) to be cyclic. Equality in the first bound forces the differential of \(P\otimes M\) to vanish, and hence \(I/I^2\) is free for \(M=R/I\). The concluding use of the Ferrand–Vasconcelos criterion requires finite projective dimension, which is available here. That cited criterion itself was not reproved.\par
\subsection*{The asymptotic inequality}
For the proof of Theorem 4.1, write \[a_j=h_j^\infty(\widetilde\Lambda^2P),\qquad
 s_j=h_j^\infty(\widetilde S^2P).\] The long exact sequence and Theorem 2.26 give \(a_j=s_j\) for \(j\ge3\). Equations (4.8)–(4.9), p. 15, then give \[\begin{aligned}
 2^d h^\infty(P)&\le a_1+s_0-a_0+s_1+a_2-s_2,\\
 h^\infty(P\otimes P)&\ge s_0+\sum_{j\ge0}|a_j-s_{j+1}|\\
 &\ge s_0-a_0+a_1+s_1+a_2-s_2.\end{aligned}\] For the last inequality, use \(s_1-a_0\) in the first absolute value, \(a_1-s_2\) in the second, and telescope the finite tail from \(j=2\). This checks the signs and the tail endpoint explicitly. Lemma 4.4 supplies \[h^\infty(P\otimes P)\le h^\infty(P)\beta(P).\] A minimal representative remains minimal under Frobenius, so its Betti sum is unchanged; the bounded spectral sequence gives the required length bound. Under Theorem 4.1’s nonzero-homology assumption, the Hilbert–Kunz argument when \(H_0(P)\ne0\), and the tiny-complex reduction when \(H_0(P)=0\), justify \(h^\infty(P)>0\).\par
For Theorem 4.10, the asymptotic argument needs only \(a_3\ge s_3\), obtained from the long exact sequence and asymptotic vanishing. The stronger finite-iterate injection mentioned on p. 16 is unnecessary for this inference. No additional defect is inferred merely from that stronger wording.\par
\section*{Assessment and limits}
The principal finding is the false short-complex domain in Definition 1.5 and its use in Theorem 5.1. A separate nonzero-homology restriction is needed in the formulations identified above. These conclusions are exact algebraic examples, not numerical evidence or an inference from missing code.\par
The full text was inspected, but the full proofs of cited background theorems were not independently reproduced. No issue was established in the checked finite-length characteristic-two inequality or in the checked equality deduction, subject to their stated external dependencies. The publisher version may differ and has not been adjudicated.\par
This is an initial scientific review. The scientific decision remains pending; there are no human reviews, no author response, and no independent acceptance. They were not separate human reviewers or independently verified alternative models. The external controller’s UI observations are recorded as such.\par
9 K. VandeBogert and M. E. Walker, The total rank conjecture in characteristic 2, arXiv:2305.09771v2 (2024), https://arxiv.org/abs/2305.09771v2. Associated journal DOI: https://doi.org/10.1215/00127094-2024-0027.\par
\end{document}
